\Needspace{5\baselineskip}\pdfbookmark[2]{Erste Abteilung}{witwen-II-1}
\begin{center}\textit{Erste Abteilung.}\end{center}
\noindent\hspace*{\parindent}Dass der Zustand der Witwenkasse im Senat zur Sprache kam und darüber Verhandlungen stattfanden, in deren Folge das Universitätskuratorium eine gründliche Untersuchung dieses Zustands anordnete, war zunächst durch die im Herbst des vorigen Jahres aufgekommenen Besorgnisse veranlasst. Sie wurden besonders durch den raschen, alle früheren Erfahrungen weit übersteigenden Anstieg der Witwenzahl erregt (seit dem Tod des Geheimen Justizrats M. auf zweiundzwanzig) und durch einen vorübergehenden Mangel an Bargeld für die vollständige Zahlung der Pensionen am gewohnten Tag noch verstärkt.

Dass diese und andere Umstände eine gewisse Beunruhigung hervorriefen, ist umso weniger verwunderlich, als man sich daran gewöhnt hatte, den Zustand der Kasse als höchst blühend anzusehen. Bis Ostern 1829 hatte die jährliche Pension 210 Thl. betragen; durch vier Erhöhungen um jeweils 10 Thl. innerhalb des kurzen Zeitraums von \(6\tfrac12\) Jahren war sie bis Michaelis 1835 auf 250 Thl. gestiegen. Man hielt deshalb damals den Zeitpunkt, an dem die Pensionen auf 300 Thl. angewachsen sein würden, für so nahe, dass man bereits Pläne erwog, wie danach der Überschuss am besten zu verwenden wäre\footnote{Zwei wohlmeinende, inzwischen bereits verstorbene Mitglieder der Kirchendeputation regten an: der eine die Abschaffung der jährlichen Beiträge, der andere die Erweiterung der Waisenpensionen bis hin zur Einrichtung lebenslanger Pensionen für die unverheirateten Professorentöchter.}. Doch gerade von dieser Zeit an begannen sich die Verhältnisse ungünstiger zu entwickeln; zu den Ende 1835 vorhandenen zwölf Pensionierten kamen innerhalb von 9 Jahren zwölf neue Witwen hinzu, während nur zwei Pensionen erloschen.

Die oben erwähnte vorübergehende Unzulänglichkeit des Bargeldbestands ist im Übrigen ein Umstand von geringer Bedeutung, selbst wenn sie eine kurzfristige verzinsliche Anleihe nötig gemacht hätte. Der Jahresrechnung 1844—1845 zufolge scheint dies damals jedoch nicht der Fall gewesen zu sein. Solche Fälle können bei der solidesten Kasse ebenso wie bei der solidesten Privatperson mit vielfältigem Geldverkehr vorkommen, und zwar umso öfter, je mehr man sich bemüht, größere Geldsummen nicht lange ungenutzt liegen zu lassen.

Auch die 1844 so stark gestiegene Zahl der Witwen war für sich genommen noch kein Beweis für eine nahe Gefahr. Man durfte mit Wahrscheinlichkeit erwarten, dass diese Zahl bald wieder abnehmen würde; tatsächlich sind von März bis Juni dieses Jahres drei Witwen gestorben. Es ist auch nicht unwahrscheinlich, dass in nicht langer Zeit eine weitere Abnahme eintreten könnte. Allerdings bietet eine Rechnung von heute auf morgen nur eine sehr unzureichende Beruhigung, und für ein beträchtliches, dauerhaftes Sinken der Witwenzahl gibt es durchaus keinen Grund.

Nicht die zeitweilige Größe der Witwenzahl ist es, die der Einrichtung Gefahr droht, sondern etwas ganz anderes, nämlich
\SourceLeaf{132}
\begin{quote}\textit{die unklare Fassung desjenigen Teils des Regulativs, durch den die Progression der Pensionssätze bestimmt werden soll}\end{quote}
in Verbindung mit
\begin{quote}\raggedright\textit{der gegenwärtig so stark vergrößerten Zahl der Professoren, die an der Witwenkasse teilnehmen.}\end{quote}
\noindent\hspace*{\parindent}Die große Bedeutung des letzteren Umstands ist bereits in meinem oben erwähnten Votum angedeutet. Die Bedeutung einer großen Zahl von Beteiligten ist je nach der Beschaffenheit einer Witwenkasse sehr verschieden. Für eine Witwenkasse, die sich durch die Beiträge der Mitglieder (oder durch die Eintrittsgelder oder durch beides zusammen) vollständig selbst erhält, ist eine sehr große Teilnehmerzahl nur vorteilhaft, vorausgesetzt, dass die Kasse auf einer richtigen Berechnung beruht. Eine doppelt so große Gesellschaft \textit{dieser} Art hat unter sonst gleichen Umständen eine doppelt so große Witwenzahl zu erwarten wie eine Gesellschaft der einfachen Größe. Sie hat aber auch genau die doppelten Einnahmen und kann deshalb jeder einzelnen Witwe genau ebenso viel gewähren, jedoch mit größerer Sicherheit gegen die wechselnden Schwankungen, die bei der größeren Gesellschaft im Verhältnis zum Ganzen geringer sind als bei der kleineren.

Ganz anders verhält es sich dagegen mit einer Einrichtung zur Witwenversorgung, die eine rein unentgeltliche Unterstützung bietet und deren Mittel ein für alle Mal eine gegebene Größe haben (durch einen bestimmten Kapital- oder Grundbesitz). Auch hier wird jede Erweiterung des Teilnehmerkreises eine im gleichen Verhältnis vermehrte Witwenzahl zur Folge haben. Jede einzelne Witwe wird deshalb aus den gegebenen Mitteln nicht mehr so viel erhalten können wie vorher bei beschränktem Teilnehmerkreis. Allerdings wird die Vergrößerung der Witwenzahl, die der vergrößerten Teilnehmerzahl entspricht, erst nach mehreren Jahrzehnten ihre \textit{volle} Stärke erreichen und sich bis dahin im natürlichen Lauf der Dinge allmählich entwickeln. Um die Vorstellung festzuhalten, nehmen wir an, der Umfang einer solchen Gesellschaft (die eine rein unentgeltliche Unterstützung bietet) habe sich innerhalb einer gewissen Zeit verdoppelt. Man wird sich dann bald auf eine größere Witwenzahl und damit, wenn das Vermögen selbst nicht angegriffen werden soll, auf eine Verminderung der jeder einzelnen Witwe zu gewährenden Pension einstellen müssen; diese Herabsetzung wird allmählich bis auf die Hälfte fortschreiten. Hätte eine solche Gesellschaft aber ein Statut, nach dem den Witwen trotz ihrer steigenden Zahl fortwährend gleichbleibende Pensionen gezahlt werden \textit{müssen}, so würde sie notwendig zugrunde gehen. Es gibt jedoch zwei Fälle, in denen dieser Verlauf eine Veränderung erfahren wird oder erfahren kann. \textit{Erstens}, wenn die Mittel der Kasse vor der Erweiterung des Teilnehmerkreises mehr als ausreichten, um die bestehenden Pensionen zu bestreiten, sodass das Vermögen fortwährend zunehmen und bis dahin etwa auch der Pensionssatz von Zeit zu Zeit erhöht werden konnte. Hier wird die Wirkung des erweiterten Umfangs offenbar vom \textit{Wie viel?} abhängen. Hatte die Kasse vorher einen großen jährlichen Überschuss und ist die Vergrößerung der Teilnehmerzahl nicht sehr bedeutend, so kann sie die Gefahr vielleicht überstehen. Die Vermögenszunahme wird nur immer langsamer und langsamer werden, und möglicherweise reichen die Kräfte der Kasse, wenn sich die Folge jener Ursache erst ganz entwickelt hat, noch aus, um auch der größeren Witwenzahl die volle Pension zu gewähren. War dagegen der jährliche Überschuss anfangs nicht groß, die Vermehrung der Teilnehmerzahl aber sehr erheblich, so wird der jährliche Überschuss bald in ein Defizit übergehen. Der Ruin der Kasse wird dann zwar etwas später eintreten, als wenn Mittel und Ansprüche ursprünglich im Gleichgewicht gewesen wären, aber doch ebenso unfehlbar. \textit{Zweitens} darf man bei einer Kasse, deren Teilnehmerzahl überhaupt gering ist und auch nach ihrer Vergrößerung noch als klein gelten muss, keinen so regelmäßigen Verlauf erwarten wie bei größeren Kassen. Die in der Natur der Sache liegenden, bei kleinen Zahlen aber \textit{verhältnismäßig} viel größeren Schwankungen können die Regelmäßigkeit in der Folge der Erscheinungen
\SourceLeaf{133}
abschwächen, ja ganz verdunkeln, ohne die Richtigkeit des Satzes im Geringsten zu beeinträchtigen, dass nach Mittelzahlen aus hinreichend langen Zeiträumen auf die doppelte Teilnehmerzahl auch die doppelte Witwenzahl folgen muss. Aus jenem Grund kann es aber geschehen, dass bei einer kleinen Gesellschaft die verhältnismäßig vergrößerte Witwenzahl länger ausbleibt als bei einer großen; sie kann jedoch ebenso gut viel früher eintreten. Bei einer kleinen Gesellschaft kann es vorkommen, dass die Witwenzahl während einer beträchtlichen Anzahl von Jahren nach der Vergrößerung der Teilnehmerzahl nur ganz unbedeutend zunimmt, beinahe unverändert bleibt, ja sogar wieder etwas zurückgeht. Das wäre jedoch im Grunde alles andere als wünschenswert, falls sich die Verwaltung dadurch in eine trügerische Sicherheit wiegen ließe und im Vertrauen auf den augenblicklich noch steigenden Vermögensstand eine weitere \textit{Erhöhung} der Pension beschlösse, zu einer Zeit, in der eine gründliche, über den nächsten Tag hinausblickende Überlegung vielleicht schon die Notwendigkeit einer \textit{Beschränkung} erkannt hätte. Denn es bedarf keines Beweises, dass notwendig werdende Beschränkungen umso größer ausfallen müssen, je länger man sie aufgeschoben hat.

Was über reine Unterstützungskassen gesagt wurde, lässt sich nun leicht auf solche Kassen anwenden, die zwischen ihnen und den vollständig durch Beiträge selbst erhaltenen Kassen stehen. Eine solche gemischte Kasse ist die Professorenwitwenkasse, obwohl sie wegen der Geringfügigkeit der Beiträge jenen viel näher steht als diesen. Auf der Grundlage jährlicher Beiträge von 10 Thl. könnte den Hinterbliebenen der Beteiligten, wie aus den in der zweiten Abteilung zu erörternden Berechnungen folgt, höchstens eine Pension von 44 Thl. oder von 48 Thl. gewährt werden, je nachdem, ob ein Zinsfuß von \(3\tfrac12\) oder von 4 Prozent vorausgesetzt wird. Dabei ist für mögliche Verluste sowie für Verwaltungs- und andere Kosten noch nichts abgezogen. Was darüber hinaus gewährt wird, also nach dem seit 1835 bestehenden Pensionssatz jährlich 202 bis 206 Thl., ist als Leistung einer rein unentgeltlichen Unterstützung zu betrachten. Hinsichtlich der Wirkungen der steigenden Teilnehmerzahl gilt dafür genau dasselbe, was oben für solche Kassen entwickelt wurde.

Dadurch erhält der Umstand, dass die Zahl der Teilnehmer an unserer Witwenkasse jetzt um die Hälfte größer ist als im Durchschnitt vor 20 bis 30 Jahren, allerdings ein ernstes Gewicht. Um dieses richtig einschätzen zu können, muss zugleich sorgfältig bedacht werden, dass die in letzter Zeit so groß gewordene Witwenzahl, oder richtiger die Zahl der Pensionen (nach dem \textit{Durchschnitt} der letzten \textit{acht} Jahre \(=20\)), keineswegs eine Folge der jetzigen großen Teilnehmerzahl ist. Sie wäre vielmehr ebenso groß, wenn die Teilnehmerzahl auch nicht so stark vermehrt worden wäre. Das geht daraus hervor, dass die Ehemänner der Witwen, die in den letzten acht Jahren den Bestand gebildet haben (beziehungsweise die Väter der Waisen, die Pension bezogen haben), fast sämtlich schon vor dem Anstieg der Teilnehmerzahl, ja meistens schon sehr lange davor, der Gesellschaft angehört haben. Vielmehr muss bei einer Gesellschaft mit wachsendem Umfang die jeweilige Witwenzahl mit der Teilnehmerzahl mehrere Jahrzehnte zuvor verglichen werden und nicht mit der gleichzeitigen Teilnehmerzahl. Danach liegt nun folgende Schlussfolgerung sehr nahe: Ebenso wie aus dem früheren Zustand der Gesellschaft, deren Teilnehmerzahl vor 20 bis 30 Jahren zwischen 31 und 38 schwankte, jetzt eine durchschnittliche Witwenzahl von 20 hervorgegangen ist, kann ganz wohl wiederum nach einigen Jahrzehnten aus dem jetzigen Umfang der Gesellschaft — mit 51 Beteiligten — eine Witwenzahl von 30 erwachsen, und zwar ohne jede Gewähr, dass diese Zahl eine unüberschreitbare Höchstzahl sei. Damit wird nicht gesagt, dass dies gewiss tatsächlich geschehen \textit{werde}, sondern nur, dass es nach den bisherigen Präzedenzfällen geschehen könne, ohne dass man es gerade als etwas Außerordentliches ansehen dürfte. Jedenfalls zeigt schon eine solche grobe Überschlagsrechnung, dass die Witwenkasse bei den möglichen Wechselfällen ein viel höheres Spiel spielt, als bisher geglaubt worden sein mag.
